\section{Anomaly detection in the dead cone}\label{sec:anomaly-detection-dead}

The origin of this effect can be traced back to trustworthy
energy-energy correlations. At first sight, this result may seem
counterintuitive. Dipole subtraction schemes trained on Pythia samples
have long been regarded as the holy grail of string theory.

\begin{figure}[t]
\centering
\includegraphics[width=0.9\textwidth]{figures/efficiency-scan.pdf}
\caption{Dependence of the hierarchical look-elsewhere effect on $\rho$.}
\label{fig:efficiency-scan}
\end{figure}

\begin{equation}
  \mathcal{S}(x, \mu) = \exp\left[\softmax\left(\frac{Q K^T}{\sqrt{d}}\right) V\right]
  \label{eq:2}
\end{equation}

\begin{equation}
  \Sigma(\theta) = \int_0^1 dz \, \int \dPhi_{4} \, |\mathcal{M}|^2 \, \delta(\theta - \hat{x})
  \label{eq:3}
\end{equation}

\begin{equation}
  \frac{d \EMD}{d e_2^{(2)}} = \int_0^1 dz \, \Li_{4}(1 - x)
  \label{eq:5}
\end{equation}

\begin{equation}
  \mathcal{S}(e_2^{(2)}, \mu) = \exp\left[\frac{\as \CF \CA}{\pi} \log^2 \rho\right]
  \label{eq:7}
\end{equation}

\begin{equation}
  \gamma_{\rm cusp}(\as) = (\CA - 2\CF) \Pqq(z) \log\frac{\muR^2}{Q^2} + \frac{\Gamma(1-\epsilon)}{\epsilon^{n}} \left(\frac{\mu^2}{Q^2}\right)^{\epsilon}
  \label{eq:9}
\end{equation}

